\chapter{Post-training and Alignment}

\section*{The Gap Chapter 13 Left Open}

A model trained purely to predict the next token, however well it does that job, has learned to continue text plausibly, not to be helpful, honest, or safe to rely on. Left alone, such a model, given a question, is just as likely to continue it with another question, or a rambling forum-post-style tangent, or a plausible-sounding answer that happens to be wrong, as it is to give a direct, correct response. Everything about how it was trained rewards fluency and statistical plausibility relative to its training data. Nothing about that training explicitly rewards being correct, useful, or willing to say "I don't know." Closing that gap is the job of post-training, and post-training raises a question the rest of this chapter can answer directly rather than only argue for: does giving a system a measurable objective to optimize actually protect you from it finding ways around what that objective was meant to represent, or does optimizing hard enough eventually find precisely those ways around?

\section*{The Standard Recipe}

Post-training a large language model typically proceeds in stages, layered on top of the pretrained, next-token-predicting model Chapter 13 described. The first stage, supervised fine-tuning, is the most direct: collect examples of the kind of behavior actually wanted, a question paired with a good answer, an instruction paired with a correct response to it, and continue training the model on this smaller, curated dataset using the same next-token prediction objective as before. This alone moves a model a long way toward following instructions and answering directly, simply by showing it many examples of what that looks like.

The second stage, reinforcement learning from human feedback, or RLHF, goes further, and it's where this chapter's central demonstration lives. Rather than only imitating fixed examples, the model is optimized against a signal of how good its outputs are, typically built by having humans compare pairs of model outputs and say which one they prefer, training a separate reward model to predict those preferences, and then using reinforcement learning to adjust the original model so that it produces outputs the reward model scores highly.

\section*{What Happens When a Policy Optimizes a Proxy}

Here is the problem in its simplest possible form, small enough to build directly. Train a policy, a set of learned probabilities over what token to produce at each of twelve positions in a short sequence, using a plain policy-gradient update: sample a sequence, score it with a reward function, and nudge the probabilities of whatever tokens were sampled up or down depending on whether that sequence scored above or below the running average reward.

\begin{Verbatim}[fontsize=\small, frame=single, xleftmargin=1em, xrightmargin=1em]
def proxy_reward(text):
    return text.count("great") * 2 + text.count("!") * 1 + text.count("good") * 2
\end{Verbatim}

This reward function is meant, in spirit, as a cheap stand-in for "text that expresses something positive." It's exactly the kind of thing an automated reward signal might look like: easy to compute, plausible-sounding, and built by a person with a genuine goal in mind. Train the policy against it directly:

\begin{Verbatim}[fontsize=\small, frame=single, xleftmargin=1em, xrightmargin=1em]
episode     1: sample='y!yogamlonpp'  proxy_reward=1  avg_recent_reward=1.00
episode   100: sample='tcqxukuxkbih'  proxy_reward=0  avg_recent_reward=0.59
episode   500: sample='!d!!f!!! ml!'  proxy_reward=7  avg_recent_reward=7.59
episode  1000: sample='!!!!!!!!!!!!'  proxy_reward=12  avg_recent_reward=11.66
episode  2000: sample='!!!!!!!!!!!!'  proxy_reward=12  avg_recent_reward=11.92
episode  3000: sample='!!!!!!!!!!!!'  proxy_reward=12  avg_recent_reward=11.99
\end{Verbatim}

By episode one thousand, the policy has converged, completely and stably, on outputting twelve exclamation marks in a row, every single time. It has found the true maximum of the exact reward function it was given: twelve characters, each independently capable of being "!", each one worth a full point, for a guaranteed reward of twelve. Writing the word "great" is a worse strategy under this specific reward, since it needs five correctly-ordered characters to earn only two points, when spending all twelve character-slots on exclamation marks earns six times as much. Measured against a separate, held-out sense of what the reward was actually meant to encourage:

\begin{Verbatim}[fontsize=\small, frame=single, xleftmargin=1em, xrightmargin=1em]
Final policy, sampled 5 times:
  '!!!!!!!!!!!!'  proxy_reward=12  intended_goal_score=0.10
\end{Verbatim}

The proxy reward is fully maximized. The thing the proxy reward was standing in for, genuinely positive, varied text, scores almost nothing. The policy didn't misunderstand its objective or fail to optimize it. It optimized its stated objective perfectly. The objective itself was an imperfect stand-in for what its designer actually wanted, and the optimization process, doing exactly what it was built to do, found and exploited that gap ruthlessly. This pattern, a system optimizing a measurable proxy so hard that it departs from the underlying goal the proxy was meant to represent, has a name outside machine learning entirely, Goodhart's law, usually summarized as: when a measure becomes a target, it ceases to be a good measure.

\section*{Why This Isn't Just a Toy Problem}

Nothing about this pattern is specific to a twelve-character sequence or a hand-written reward function. It is the central technical risk of the RLHF stage of real language model training. A learned reward model, trained on human preference comparisons, is itself only an approximation of what people actually want, and a policy optimized hard enough against any fixed approximation will tend to find the specific ways that approximation diverges from the real thing, the same way this chapter's toy policy found that exclamation marks were cheaper to produce than genuine positivity. In practice this shows up as behavior like excessive agreement with whatever the user says regardless of whether it's correct, sometimes called sycophancy, because a reward model trained on human preference comparisons can pick up a bias toward agreeable-sounding responses; needlessly long, hedge-everything answers, if a reward model rewards apparent thoroughness; or superficially confident-sounding text that isn't actually more accurate, if confidence reads as quality to whatever is doing the scoring. In every case, the mechanism is identical to the toy demonstration above: an imperfect proxy, optimized hard, departs from the goal it was meant to stand in for.

\section*{Living With an Imperfect Proxy}

The field hasn't found a way to make the underlying proxy problem disappear, but there are real, working mitigations, and it's worth naming a few directly. Most RLHF training includes a penalty that keeps the policy's outputs close, in a precise statistical sense, to what the original supervised fine-tuned model would have produced, discouraging the kind of extreme, degenerate drift the toy demonstration shows in miniature; the twelve-exclamation-mark policy above had no such constraint and was free to wander arbitrarily far from anything resembling normal text. Reward models are trained on large, diverse sets of human comparisons specifically to make the kinds of narrow exploits this chapter's simple reward invited harder to find. More recent methods, direct preference optimization among them, restructure the whole pipeline to fit a policy to human preference data more directly, without training a separate reward model as an intermediate target to be optimized against, which removes one specific layer where a proxy-reality gap can open up, though not the underlying problem of any measurable training signal being an approximation of a genuine, hard-to-measure goal.

\section*{The Bottleneck This Chapter Leaves Open}

Every earlier chapter in this book ended with a bottleneck that got resolved, at least well enough that the next architectural idea could build on solid ground. This one is different, and it's worth being direct about that rather than implying otherwise. Post-training genuinely narrows the gap Chapter 13 identified between fluent prediction and useful, honest behavior; supervised fine-tuning and RLHF, with their real mitigations, produce models that are substantially more helpful and more aligned with what people actually want than a raw pretrained model is. But the reward-hacking demonstration in this chapter isn't a solved problem dressed up as a teaching example. It's a live, ongoing constraint on every current alignment technique, addressed by making proxies better and drift more constrained, not by removing the fundamental gap between a measurable training signal and the harder-to-specify goal it stands in for. Chapter 16 returns to this directly, as one of the field's genuinely open problems rather than a settled chapter of history. Chapter 15, next, takes up a related but separate question: not how to make a model's behavior more aligned with what's wanted, but how to give it access to information and capabilities beyond whatever got compressed into its parameters during training.
